\documentclass[11pt]{article}
\usepackage{mathblatt}


\begin{document}
\pfheftstil
\pfheftkopf{Quadratische}{P10}
\pfrueckblick
\begin{pfaufg}{}{1.}{}
Schreib als x $-$ d und gib d an: x + 3; x $-$ 2; x + 1.
\par\smallskip\noindent \begingroup\renewcommand{\arraystretch}{1.3}\begin{tabular}{|c|c|}\hline \textbf{Term} & \textbf{d} \\ \hline x $-$ ($-$3), d & \rule{0pt}{3ex}\hspace{1.6cm} \\ \hline x $-$ 2, d & \rule{0pt}{3ex}\hspace{1.6cm} \\ \hline x $-$ ($-$1), d & \rule{0pt}{3ex}\hspace{1.6cm} \\ \hline\end{tabular}\endgroup\par
\fusshilfe{\mbox{1:~$-$3; 2; $-$1}}
\end{pfaufg}
\begin{pfaufg}{}{2.}{}
Multipliziere aus: (x $-$ 3)²; (x + 3)²; (x $-$ 2)².
\rechenplatz{1}
\fusshilfe{\mbox{2:~x² $-$ 6x + 9; x² + 6x + 9; x² $-$ 4x + 4}}
\end{pfaufg}
\begin{pfaufg}{}{3.}{}
Ergänze die Tabelle Gleichung | p | q | (p/2)² $-$ q.
\par\smallskip\noindent \begingroup\renewcommand{\arraystretch}{1.3}\begin{tabular}{|c|c|c|c|}\hline \textbf{Gleichung} & \textbf{p} & \textbf{q} & \textbf{(p/2)² $-$ q} \\ \hline $-$6 | 7 | 2 & \rule{0pt}{3ex}\hspace{1.6cm} & \rule{0pt}{3ex}\hspace{1.6cm} & \rule{0pt}{3ex}\hspace{1.6cm} \\ \hline 4 | 3 | 1 & \rule{0pt}{3ex}\hspace{1.6cm} & \rule{0pt}{3ex}\hspace{1.6cm} & \rule{0pt}{3ex}\hspace{1.6cm} \\ \hline $-$2 | $-$3 | 4 & \rule{0pt}{3ex}\hspace{1.6cm} & \rule{0pt}{3ex}\hspace{1.6cm} & \rule{0pt}{3ex}\hspace{1.6cm} \\ \hline\end{tabular}\endgroup\par
\fusshilfe{\mbox{3:~$-$6 | 7 | 2; 4 | 3 | 1; $-$2 | $-$3 | 4}}
\end{pfaufg}
\begin{pfaufg}{}{4.}{}
Bring alles auf eine Seite: x² + 3x = 2x + 6; x² $-$ 1 = $-$2x + 2.
\rechenplatz{1}
\fusshilfe{\mbox{4:~x² + x $-$ 6 = 0; x² + 2x $-$ 3 = 0}}
\end{pfaufg}
\pfstufekopf[sscheitelpunktablesen]{Scheitelpunkt ablesen}{in 5 der letzten 5 Prüfungen}
\pfgruppe{}
\begin{pfaufg}{}{5.}{}
Kreise in $f(x) = (x - 4)^2 + 2$ die Zahl in der Klammer ein. Bei welchem $x$ liegt der Scheitel?
\par\noindent x = \leerfeld \par
\fusshilfe{\mbox{5:~$4$}}
\end{pfaufg}
\begin{pfaufg}{}{6.}{}
Kreuze an, in welcher Form $f(x) = x^2 + 4x - 1$ geschrieben ist.
\pfkreuzzeile{Scheitelpunktform}\pfkreuzzeile{Normalform}
\fusshilfe{\mbox{6:~Normalform}}
\end{pfaufg}
\pfgruppe{}
\begin{pfaufg}{P10 ’17}{7.}{}
Die Parabel $p$ hat die Gleichung $y = (x+3)^2 - 2$ und ist im Koordinatensystem gezeichnet. Gib den Scheitelpunkt der Parabel $p$ an: $S(\ \ |\ \ )$.
\par\smallskip\noindent \begin{tikzpicture}\begin{axis}[xmin=-6,xmax=2,ymin=-2,ymax=4,axis lines=middle,tick label style={font=\scriptsize},clip=true,/pgf/number format/use comma,axis line style={-{Stealth[length=2mm]}},every axis plot/.append style={line width=0.9pt},x=0.55cm,y=0.55cm,xtick distance=1,ytick distance=1,enlargelimits=false,grid=both,grid style={mbgitter,line width=0.3pt},major grid style={mbgitter},xlabel={$x$},ylabel={$y$},x label style={anchor=west},y label style={anchor=south}]\addplot[black,domain=-6:2,samples=120] {((x+3)^2)-2};\node[font=\small,fill=white,inner sep=1pt,anchor=south west] at (axis cs:-0.752,3.0535) {$p$};\end{axis}\end{tikzpicture}\par
\rechenplatz{1}
\fusshilfe{\mbox{7:~S($-$3 | $-$2)}}
\end{pfaufg}
\begin{pfaufg}{P10 ’18}{8.}{}
Im Koordinatensystem ist die Parabel $p$ mit $p(x) = x^2 - 4x + 2$ gezeichnet. Gib den Scheitelpunkt von $p$ an: $S(\ \ |\ \ )$.
\par\smallskip\noindent \begin{tikzpicture}\begin{axis}[xmin=-4,xmax=5,ymin=-2,ymax=7,axis lines=middle,tick label style={font=\scriptsize},clip=true,/pgf/number format/use comma,axis line style={-{Stealth[length=2mm]}},every axis plot/.append style={line width=0.9pt},x=0.55cm,y=0.55cm,xtick distance=1,ytick distance=1,enlargelimits=false,grid=both,grid style={mbgitter,line width=0.3pt},major grid style={mbgitter},xlabel={$x$},ylabel={$y$},x label style={anchor=west},y label style={anchor=south}]\addplot[black,domain=-4:5,samples=120] {(x^2)-4*x+2};\node[font=\small,fill=white,inner sep=1pt,anchor=south west] at (axis cs:4.28,3.1984) {$p$};\end{axis}\end{tikzpicture}\par
\rechenplatz{1}
\fusshilfe{\mbox{8:~S(2 | $-$2)}}
\end{pfaufg}
\begin{pfaufg}{P10 ’20}{9.}{}
Die Parabel zu $f$ mit $y = (x+2)^2 - 4$ ist im Koordinatensystem gezeichnet. $S$ ist der Scheitelpunkt des Graphen von $f$. Gib seine Koordinaten an.
\par\smallskip\noindent \begin{tikzpicture}\begin{axis}[xmin=-5,xmax=7,ymin=-5,ymax=5,axis lines=middle,tick label style={font=\scriptsize},clip=true,/pgf/number format/use comma,axis line style={-{Stealth[length=2mm]}},every axis plot/.append style={line width=0.9pt},x=0.55cm,y=0.55cm,xtick distance=1,ytick distance=1,enlargelimits=false,grid=both,grid style={mbgitter,line width=0.3pt},major grid style={mbgitter},xlabel={$x$},ylabel={$y$},x label style={anchor=west},y label style={anchor=south}]\addplot[black,domain=-5:7,samples=120] {((x+2)^2)-4};\node[font=\small,fill=white,inner sep=1pt,anchor=south west] at (axis cs:0.76,3.6176) {$f$};\end{axis}\end{tikzpicture}\par
\rechenplatz{1}
\fusshilfe{\mbox{9:~S($-$2 | $-$4)}}
\end{pfaufg}
\begin{pfaufg}{P10 ’22}{10.}{}
Die lineare Funktion $f$ hat die Gleichung $f(x) = -2x + 3$. Die verschobene Normalparabel $p$ hat die Gleichung $p(x) = (x-2)^2 - 4$. Gib den Scheitelpunkt $S$ von $p$ an.
\rechenplatz{1}
\fusshilfe{\mbox{10:~S(2|$-$4)}}
\end{pfaufg}
\pfgruppe{}
\begin{pfaufg}{P10 ’21}{11.}{}
Kreuze den Scheitelpunkt der gezeichneten Parabel an:
\pfkreuzzeile{$S(-1\,|\,-2)$}\pfkreuzzeile{$S(-1\,|\,2)$}\pfkreuzzeile{$S(-2\,|\,-1)$}\pfkreuzzeile{$S(2\,|\,-1)$}
\par\smallskip\noindent \begin{tikzpicture}\begin{axis}[xmin=-4,xmax=2,ymin=-2,ymax=3,axis lines=middle,tick label style={font=\scriptsize},clip=true,/pgf/number format/use comma,axis line style={-{Stealth[length=2mm]}},every axis plot/.append style={line width=0.9pt},x=0.55cm,y=0.55cm,xtick distance=1,ytick distance=1,enlargelimits=false,grid=both,grid style={mbgitter,line width=0.3pt},major grid style={mbgitter},xlabel={$x$},ylabel={$y$},x label style={anchor=west},y label style={anchor=south}]\addplot[black,domain=-4:2,samples=120] {((x+2)^2)-1};\end{axis}\end{tikzpicture}\par
\fusshilfe{\mbox{11:~S($-$2 | $-$1)}}
\end{pfaufg}
\pfsteckt{steckt auch in Nr. 18 (P10 ’23); Nr. 39 (P10 ’24); Nr. 17 (P10 ’25)}
\pfstufekopf[sscheitelpunktformang]{Scheitelpunktform angeben}{in 2 der letzten 5 Prüfungen}
\pfgruppe{}
\begin{pfaufg}{P10 ’16}{12.}{}
Der Scheitelpunkt einer verschobenen Normalparabel ist $S(1\,|\,3)$. Kreuze die Gleichung an, die passt:
\pfkreuzzeile{$y = (x+1)^2 + 3$}\pfkreuzzeile{$y = (x-1)^2 - 3$}\pfkreuzzeile{$y = (x-1)^2 + 3$}
\fusshilfe{\mbox{12:~y = (x $-$ 1)² + 3}}
\end{pfaufg}
\pfgruppe{}
\begin{pfaufg}{P10 ’15}{\llap{\textasteriskcentered\,}13.}{}
Die gezeichnete Normalparabel wird um 2 Einheiten nach rechts verschoben. Gib den Scheitelpunkt der verschobenen Parabel an: $S(\ \ |\ \ )$.
\par\smallskip\noindent \begin{tikzpicture}\begin{axis}[xmin=-2,xmax=2,ymin=-1,ymax=4,axis lines=middle,tick label style={font=\scriptsize},clip=true,/pgf/number format/use comma,axis line style={-{Stealth[length=2mm]}},every axis plot/.append style={line width=0.9pt},x=0.55cm,y=0.55cm,xtick distance=1,ytick distance=1,enlargelimits=false,grid=both,grid style={mbgitter,line width=0.3pt},major grid style={mbgitter},xlabel={$x$},ylabel={$y$},x label style={anchor=west},y label style={anchor=south}]\addplot[black,domain=-2:2,samples=120] {(x^2)};\end{axis}\end{tikzpicture}\par
\rechenplatz{1}
\fusshilfe{\mbox{13:~S(2 | 0)}}
\end{pfaufg}
\begin{pfaufg}{P10 ’18}{\llap{\textasteriskcentered\,}14.}{}
Im Koordinatensystem ist die Parabel $p$ mit $p(x) = x^2 - 4x + 2$ gezeichnet. Schreib die Gleichung von $p$ in Scheitelpunktform.
\par\smallskip\noindent \begin{tikzpicture}\begin{axis}[xmin=-4,xmax=5,ymin=-2,ymax=7,axis lines=middle,tick label style={font=\scriptsize},clip=true,/pgf/number format/use comma,axis line style={-{Stealth[length=2mm]}},every axis plot/.append style={line width=0.9pt},x=0.55cm,y=0.55cm,xtick distance=1,ytick distance=1,enlargelimits=false,grid=both,grid style={mbgitter,line width=0.3pt},major grid style={mbgitter},xlabel={$x$},ylabel={$y$},x label style={anchor=west},y label style={anchor=south}]\addplot[black,domain=-4:5,samples=120] {(x^2)-4*x+2};\node[font=\small,fill=white,inner sep=1pt,anchor=south west] at (axis cs:4.28,3.1984) {$p$};\end{axis}\end{tikzpicture}\par
\rechenplatz{1}
\fusshilfe{\mbox{14:~p(x) = (x $-$ 2)² $-$ 2}}
\end{pfaufg}
\begin{pfaufg}{P10 ’20}{\llap{\textasteriskcentered\,}15.}{}
Die Parabel zu $f$ mit $y = (x+2)^2 - 4$ ist im Koordinatensystem gezeichnet. Spiegelt man den Graphen von $f$ an der y-Achse, entsteht die Parabel $p$. Skizziere $p$ in das Koordinatensystem. Gib eine Gleichung von $p$ an.
\par\smallskip\noindent \begin{tikzpicture}\begin{axis}[xmin=-5,xmax=7,ymin=-5,ymax=5,axis lines=middle,tick label style={font=\scriptsize},clip=true,/pgf/number format/use comma,axis line style={-{Stealth[length=2mm]}},every axis plot/.append style={line width=0.9pt},x=0.55cm,y=0.55cm,xtick distance=1,ytick distance=1,enlargelimits=false,grid=both,grid style={mbgitter,line width=0.3pt},major grid style={mbgitter},xlabel={$x$},ylabel={$y$},x label style={anchor=west},y label style={anchor=south}]\addplot[black,domain=-5:7,samples=120] {((x+2)^2)-4};\node[font=\small,fill=white,inner sep=1pt,anchor=south west] at (axis cs:0.76,3.6176) {$f$};\end{axis}\end{tikzpicture}\par
\rechenplatz{1}
\fusshilfe{\mbox{15:~y = (x $-$ 2)² $-$ 4}}
\end{pfaufg}
\begin{pfaufg}{P10 ’17}{\llap{\textasteriskcentered\,}16.}{}
Die Parabel $p$ hat die Gleichung $y = (x+3)^2 - 2$ und ist im Koordinatensystem gezeichnet. Man verschiebt $p$ um 2 Einheiten nach oben und spiegelt das Ergebnis danach an der x-Achse. So entsteht die Parabel $q$. Gib eine Gleichung von $q$ an.
\par\smallskip\noindent \begin{tikzpicture}\begin{axis}[xmin=-6,xmax=2,ymin=-2,ymax=4,axis lines=middle,tick label style={font=\scriptsize},clip=true,/pgf/number format/use comma,axis line style={-{Stealth[length=2mm]}},every axis plot/.append style={line width=0.9pt},x=0.55cm,y=0.55cm,xtick distance=1,ytick distance=1,enlargelimits=false,grid=both,grid style={mbgitter,line width=0.3pt},major grid style={mbgitter},xlabel={$x$},ylabel={$y$},x label style={anchor=west},y label style={anchor=south}]\addplot[black,domain=-6:2,samples=120] {((x+3)^2)-2};\node[font=\small,fill=white,inner sep=1pt,anchor=south west] at (axis cs:-0.752,3.0535) {$p$};\end{axis}\end{tikzpicture}\par
\rechenplatz{2}
\fusshilfe{\mbox{16:~q: y = $-$(x + 3)²}}
\end{pfaufg}
\begin{pfaufg}{P10 ’25}{17.}{}
Die quadratische Funktion $p$ hat die Gleichung $p(x) = x^2 - 6x + 7$. Ihr Graph ist im Koordinatensystem gezeichnet. Gib den Scheitelpunkt $S$ von $p$ an: $S(\ \ |\ \ )$. Schreib die Gleichung von $p$ in Scheitelpunktform: $p(x) = \ldots$
\par\smallskip\noindent \begin{tikzpicture}\begin{axis}[xmin=-3,xmax=6,ymin=-2,ymax=4,axis lines=middle,tick label style={font=\scriptsize},clip=true,/pgf/number format/use comma,axis line style={-{Stealth[length=2mm]}},every axis plot/.append style={line width=0.9pt},x=0.55cm,y=0.55cm,xtick distance=1,ytick distance=1,enlargelimits=false,grid=both,grid style={mbgitter,line width=0.3pt},major grid style={mbgitter},xlabel={$x$},ylabel={$y$},x label style={anchor=west},y label style={anchor=south}]\addplot[black,domain=-3:6,samples=120] {(x^2)-6*x+7};\node[font=\small,fill=white,inner sep=1pt,anchor=south west] at (axis cs:5.28,3.1984) {$p$};\end{axis}\end{tikzpicture}\par
\rechenplatz{2}
\fusshilfe{\mbox{17:~S(3|$-$2); p(x) = (x $-$ 3)² $-$ 2}}
\end{pfaufg}
\pfgruppe{}
\begin{pfaufg}{P10 ’23}{\llap{\textasteriskcentered\,}18.}{}
Gib den Scheitelpunkt der gezeichneten Normalparabel $p$ an: $S(\ \ |\ \ )$. Schreib eine Gleichung von $p$ in Scheitelpunktform.
\par\smallskip\noindent \begin{tikzpicture}\begin{axis}[xmin=-5,xmax=5,ymin=-4,ymax=6,axis lines=middle,tick label style={font=\scriptsize},clip=true,/pgf/number format/use comma,axis line style={-{Stealth[length=2mm]}},every axis plot/.append style={line width=0.9pt},x=0.55cm,y=0.55cm,xtick distance=1,ytick distance=1,enlargelimits=false,grid=both,grid style={mbgitter,line width=0.3pt},major grid style={mbgitter},xlabel={$x$},ylabel={$y$},x label style={anchor=west},y label style={anchor=south}]\addplot[black,domain=-5:5,samples=120] {-((x+1)^2)+6};\node[font=\small,fill=white,inner sep=1pt,anchor=south west] at (axis cs:2,-3) {$p$};\end{axis}\end{tikzpicture}\par
\rechenplatz{1}
\fusshilfe{\mbox{18:~S($-$1|6); p(x) = $-$(x + 1)² + 6}}
\end{pfaufg}
\pfgruppe{}
\begin{pfaufg}{P10 ’18}{\llap{\textasteriskcentered\,}19.}{}
Eine Normalparabel $q$ hat ihren Scheitelpunkt auf der y-Achse und keine Nullstelle. Skizziere eine solche Parabel in das Koordinatensystem und gib ihre Gleichung an. Spiegelt man $q$ an der x-Achse, entsteht $q^{\prime}$. Gib eine Gleichung von $q^{\prime}$ an.
\par\smallskip\noindent \begin{tikzpicture}\begin{axis}[xmin=-4,xmax=5,ymin=-6,ymax=6,axis lines=middle,tick label style={font=\scriptsize},clip=true,/pgf/number format/use comma,axis line style={-{Stealth[length=2mm]}},every axis plot/.append style={line width=0.9pt},x=0.55cm,y=0.55cm,xtick distance=1,ytick distance=1,enlargelimits=false,grid=both,grid style={mbgitter,line width=0.3pt},major grid style={mbgitter},xlabel={$x$},ylabel={$y$},x label style={anchor=west},y label style={anchor=south}]\end{axis}\end{tikzpicture}\par
\rechenplatz{2}
\fusshilfe{\mbox{19:~x² + 1; $-$x² $-$ 1}}
\end{pfaufg}

\pfstufekopf[snullstellenberechnen]{Nullstellen berechnen}{in 4 der letzten 5 Prüfungen}
\pfgruppe{}
\begin{pfaufg}{}{20.}{}
Die Gleichung $x^2 - x - 12 = 0$ hat die Form $x^2 + px + q = 0$. Gib $p$ und $q$ an.
\par\noindent p = \leerfeld , q = \leerfeld \par
\fusshilfe{\mbox{20:~$-12$}}
\end{pfaufg}
\begin{pfaufg}{}{21.}{}
Du willst die Schnittpunkte der Parabel $f(x) = x^2 - 2x$ und der Geraden $g(x) = x + 4$ berechnen. Schreibe nur die Gleichung auf, mit der du anfängst.
\par\noindent \leerfeld[=]  \leerfeld \par
\fusshilfe{\mbox{21:~$x + 4$}}
\end{pfaufg}
\pfgruppe{}
\begin{pfaufg}{NRW ’24}{22.}{}
Löse die Gleichung $-$5x² + 20 = 0.
\rechenplatz{2}
\fusshilfe{\mbox{22:~x² = 4 $\Rightarrow$ x$_1$ = $-$2, x$_2$ = 2}}
\end{pfaufg}
\begin{pfaufg}{P10 ’20}{\llap{\textasteriskcentered\,}23.}{}
Berechne die Nullstellen der Funktion mit $y = 2x^2 + 8x + 6$.
\rechenplatz{3}
\fusshilfe{\mbox{23:~x$_1$ = $-$1; x$_2$ = $-$3}}
\end{pfaufg}
\pfgruppe{}
\begin{pfaufg}{P10 ’25}{\llap{\textasteriskcentered\,}24.}{}
Die quadratische Funktion $p$ hat die Gleichung $p(x) = x^2 - 6x + 7$. Ihr Graph ist im Koordinatensystem gezeichnet. Berechne die Koordinaten der Punkte, in denen die Parabel $p$ die x-Achse schneidet.
\par\smallskip\noindent \begin{tikzpicture}\begin{axis}[xmin=-3,xmax=6,ymin=-2,ymax=4,axis lines=middle,tick label style={font=\scriptsize},clip=true,/pgf/number format/use comma,axis line style={-{Stealth[length=2mm]}},every axis plot/.append style={line width=0.9pt},x=0.55cm,y=0.55cm,xtick distance=1,ytick distance=1,enlargelimits=false,grid=both,grid style={mbgitter,line width=0.3pt},major grid style={mbgitter},xlabel={$x$},ylabel={$y$},x label style={anchor=west},y label style={anchor=south}]\addplot[black,domain=-3:6,samples=120] {(x^2)-6*x+7};\node[font=\small,fill=white,inner sep=1pt,anchor=south west] at (axis cs:5.28,3.1984) {$p$};\end{axis}\end{tikzpicture}\par
\rechenplatz{3}
\fusshilfe{\mbox{24:~N1(1,59|0), N2(4,41|0), x = 3 ± √2}}
\end{pfaufg}
\pfgruppe{}
\begin{pfaufg}{P10 ’17}{\llap{\textasteriskcentered\,}25.}{}
Die Parabel $p$ hat die Gleichung $y = (x+3)^2 - 2$ und ist im Koordinatensystem gezeichnet. Zeige, dass man die Gleichung von $p$ auch als $y = x^2 + 6x + 7$ schreiben kann. Berechne die Nullstellen von $p$.
\par\smallskip\noindent \begin{tikzpicture}\begin{axis}[xmin=-6,xmax=2,ymin=-2,ymax=4,axis lines=middle,tick label style={font=\scriptsize},clip=true,/pgf/number format/use comma,axis line style={-{Stealth[length=2mm]}},every axis plot/.append style={line width=0.9pt},x=0.55cm,y=0.55cm,xtick distance=1,ytick distance=1,enlargelimits=false,grid=both,grid style={mbgitter,line width=0.3pt},major grid style={mbgitter},xlabel={$x$},ylabel={$y$},x label style={anchor=west},y label style={anchor=south}]\addplot[black,domain=-6:2,samples=120] {((x+3)^2)-2};\node[font=\small,fill=white,inner sep=1pt,anchor=south west] at (axis cs:-0.752,3.0535) {$p$};\end{axis}\end{tikzpicture}\par
\rechenplatz{4}
\fusshilfe{\mbox{25:~x² + 6x + 7; $-$3 $-$ √2 $\approx$ $-$4,41}}
\end{pfaufg}
\pfsteckt{steckt auch in Nr. 28 (P10 ’21); Nr. 29 (P10 ’22); Nr. 45 (P10 ’23); Nr. 30 (P10 ’24)}
\pfstufekopf[sgeradeundparabelglei]{Gerade und Parabel gleichsetzen}{in 2 der letzten 5 Prüfungen}
\pfgruppe{}
\begin{pfaufg}{}{26.}{}
Kreuze an, welche Form die Gleichung $x^2 - 3x + 1 = 0$ hat und welcher Lösungsweg passt.
\pfkreuzzeile{reinquadratisch – Wurzelziehen}\pfkreuzzeile{Produkt gleich null – Nullprodukt}\pfkreuzzeile{quadrierte Klammer – Rückwärtsrechnen}\pfkreuzzeile{Normalform – Formel}
\fusshilfe{\mbox{26:~Formel}}
\end{pfaufg}
\pfgruppe{}
\begin{pfaufg}{nach BW ’22}{27.}{}
Die Parabeln y = x² $-$ 8x + 12 und y = (x $-$ 1)² $-$ 7 schneiden sich. Berechne den Schnittpunkt.
\rechenplatz{2}
\fusshilfe{\mbox{27:~Q(3; $-$3)}}
\end{pfaufg}
\pfgruppe{}
\begin{pfaufg}{P10 ’21}{\llap{\textasteriskcentered\,}28.}{}
Die Parabel mit $y = x^2 + 2x - 1$ und die Gerade mit $y = 3x + 1$ schneiden sich. Berechne die Koordinaten der Schnittpunkte.
\rechenplatz{4}
\fusshilfe{\mbox{28:~S$_1$($-$1 | $-$2); S$_2$(2 | 7)}}
\end{pfaufg}
\begin{pfbuendel}{gleiche Art \textperiodcentered{} 2$\times$ geprüft – kannst du überspringen}
\pfbpaar{\pfbglied{P10 ’22}{\llap{\textasteriskcentered\,}29.}{Die lineare Funktion $f$ hat die Gleichung $f(x) = -2x + 3$. Die verschobene Normalparabel $p$ hat die Gleichung $p(x) = (x-2)^2 - 4$. Bestimme die Koordinaten der Punkte, in denen die Gerade $f$ die Parabel $p$ schneidet.}{}\fusshilfe{\mbox{29:~S$_1$($-$1|5), S$_2$(3|$-$3)}}}{}
\end{pfbuendel}
\pfgruppe{}
\begin{pfaufg}{P10 ’24}{\llap{\textasteriskcentered\,}30.}{}
Die Gerade $f$ ist der Graph von $f(x) = 4x + 1$. Die Parabel $p$ ist der Graph von $p(x) = x^2 - 4$. Die Graphen von $f$ und $p$ haben zwei Schnittpunkte. Berechne ihre Koordinaten.
\rechenplatz{4}
\fusshilfe{\mbox{30:~S1($-$1|$-$3), S2(5|21)}}
\end{pfaufg}
\pfstufekopf[swertetabellezuordnen]{Wertetabelle zuordnen}{in 1 der letzten 5 Prüfungen}
\pfgruppe{}
\begin{pfaufg}{}{31.}{}
Kreuze an, ob der Graph von $f(x) = 4x - 1$ eine Gerade oder eine Parabel ist.
\pfkreuzzeile{Gerade}\pfkreuzzeile{Parabel}
\fusshilfe{\mbox{31:~Gerade}}
\end{pfaufg}
\pfgruppe{}
\begin{pfaufg}{P10 ’26}{32.}{}
Kreuze die Wertetabelle an, die zu $y = 3x^2$ gehört.
\par\smallskip\noindent \begin{tabular}[t]{@{}c@{}}\renewcommand{\arraystretch}{1.3}\begin{tabular}{|c|c|c|c|c|c|}\hline $x$ & $-$2 & $-$1 & 0 & 1 & 2 \\ \hline $y$ & 12 & 3 & 0 & 3 & 12 \\ \hline\end{tabular}\\[3pt]{\small Tabelle 1\enspace$\square$}\end{tabular}\hspace{0.5cm}\begin{tabular}[t]{@{}c@{}}\renewcommand{\arraystretch}{1.3}\begin{tabular}{|c|c|c|c|c|c|}\hline $x$ & $-$2 & $-$1 & 0 & 1 & 2 \\ \hline $y$ & $-$6 & $-$3 & 0 & 3 & 6 \\ \hline\end{tabular}\\[3pt]{\small Tabelle 2\enspace$\square$}\end{tabular}\hspace{0.5cm}\begin{tabular}[t]{@{}c@{}}\renewcommand{\arraystretch}{1.3}\begin{tabular}{|c|c|c|c|c|c|}\hline $x$ & $-$2 & $-$1 & 0 & 1 & 2 \\ \hline $y$ & 4 & 1 & 0 & 1 & 4 \\ \hline\end{tabular}\\[3pt]{\small Tabelle 3\enspace$\square$}\end{tabular}\par
\fusshilfe{\mbox{32:~erste Tabelle (y = 12, 3, 0, 3, 12)}}
\end{pfaufg}
\pfgruppe{}
\begin{pfaufg}{P10 ’15}{\llap{\textasteriskcentered\,}33.}{}
Lena verlässt ein Flugzeug mit dem Fallschirm. Die ersten 20 Sekunden stürzt sie frei, der Schirm ist noch zu; ihre Höhe sinkt dabei von 2400 m auf 1200 m. Der Graph zeigt ihre Höhe $h$ (in Metern) zur Zeit $t$ (in Sekunden). Welche Gleichung kann den Graphen in den ersten 20 Sekunden beschreiben? Kreuze an und begründe:
\pfkreuzzeile{$h(t) = 3t^2 + 2400$}\pfkreuzzeile{$h(t) = 3t^2 - 2400$}\pfkreuzzeile{$h(t) = -3t^2 + 2400$}\pfkreuzzeile{$h(t) = -3t^2 - 2400$}
\par\smallskip\noindent \begin{tikzpicture}\begin{axis}[xmin=0,xmax=280,ymin=0,ymax=2400,tick label style={font=\scriptsize},clip=true,/pgf/number format/use comma,axis line style={-{Stealth[length=2mm]}},every axis plot/.append style={line width=0.9pt},width=11.5cm,height=7.5cm,grid=both,grid style={mbgitter,line width=0.3pt},minor tick num=1,axis lines=left,xlabel={Zeit t in Sekunden},ylabel={Höhe h in Metern},label style={font=\small},scaled ticks=false,/pgf/number format/1000 sep={\,}]\addplot[black,domain=0:20,samples=120] {-3*(x^2)+2400};\node[font=\small,fill=white,inner sep=1pt,anchor=south west] at (axis cs:18.4,1384.32) {$G$};\addplot[black,domain=20:260,samples=120] {-5*x+1300};\end{axis}\end{tikzpicture}\par
\rechenplatz{1}
\fusshilfe{\mbox{33:~$-$3t² + 2400}}
\end{pfaufg}
\pfstufekopf[spunktaufderparabelpr]{Punkt auf der Parabel prüfen}{in 1 der letzten 5 Prüfungen}
\pfgruppe{}
\begin{pfaufg}{P10 ’24}{34.}{}
Die Gerade $f$ ist der Graph von $f(x) = 4x + 1$. Die Parabel $p$ ist der Graph von $p(x) = x^2 - 4$. Prüfe mit einer Rechnung, ob der Punkt $T(-5\,|\,20)$ auf der Parabel $p$ liegt.
\rechenplatz{1}
\fusshilfe{\mbox{34:~Nein, p($-$5) = 25 $-$ 4 = 21 $\neq$ 20}}
\end{pfaufg}
\pfgruppe{}
\begin{pfaufg}{P10 ’20}{35.}{}
Die Parabel zu $f$ mit $y = (x+2)^2 - 4$ ist im Koordinatensystem gezeichnet. Genau einer dieser Punkte liegt nicht auf dem Graphen von $f$. Kreuze ihn an:
\pfkreuzzeile{$A(-3\,|\,-3)$}\pfkreuzzeile{$B(0\,|\,0)$}\pfkreuzzeile{$C(-1\,|\,-3)$}\pfkreuzzeile{$D(-1{,}5\,|\,-2{,}5)$}\pfkreuzzeile{$E(-2\,|\,-4)$}
\par\smallskip\noindent \begin{tikzpicture}\begin{axis}[xmin=-5,xmax=7,ymin=-5,ymax=5,axis lines=middle,tick label style={font=\scriptsize},clip=true,/pgf/number format/use comma,axis line style={-{Stealth[length=2mm]}},every axis plot/.append style={line width=0.9pt},x=0.55cm,y=0.55cm,xtick distance=1,ytick distance=1,enlargelimits=false,grid=both,grid style={mbgitter,line width=0.3pt},major grid style={mbgitter},xlabel={$x$},ylabel={$y$},x label style={anchor=west},y label style={anchor=south}]\addplot[black,domain=-5:7,samples=120] {((x+2)^2)-4};\node[font=\small,fill=white,inner sep=1pt,anchor=south west] at (axis cs:0.76,3.6176) {$f$};\end{axis}\end{tikzpicture}\par
\fusshilfe{\mbox{35:~D($-$1,5 | $-$2,5)}}
\end{pfaufg}
\pfgruppe{}
\begin{pfaufg}{P10 ’20}{36.}{}
Die Parabel zu $f$ mit $y = (x+2)^2 - 4$ ist im Koordinatensystem gezeichnet. Der Punkt $A(2\,|\,y)$ liegt auf dem Graphen von $f$ mit $y = (x+2)^2 - 4$. Berechne $y$.
\par\smallskip\noindent \begin{tikzpicture}\begin{axis}[xmin=-5,xmax=7,ymin=-5,ymax=5,axis lines=middle,tick label style={font=\scriptsize},clip=true,/pgf/number format/use comma,axis line style={-{Stealth[length=2mm]}},every axis plot/.append style={line width=0.9pt},x=0.55cm,y=0.55cm,xtick distance=1,ytick distance=1,enlargelimits=false,grid=both,grid style={mbgitter,line width=0.3pt},major grid style={mbgitter},xlabel={$x$},ylabel={$y$},x label style={anchor=west},y label style={anchor=south}]\addplot[black,domain=-5:7,samples=120] {((x+2)^2)-4};\node[font=\small,fill=white,inner sep=1pt,anchor=south west] at (axis cs:0.76,3.6176) {$f$};\end{axis}\end{tikzpicture}\par
\rechenplatz{2}
\fusshilfe{\mbox{36:~y = 12}}
\end{pfaufg}
\pfgruppe{}
\begin{pfaufg}{P10 ’18}{37.}{}
Im Koordinatensystem ist die Parabel $p$ mit $p(x) = x^2 - 4x + 2$ gezeichnet. Entscheide für jede Aussage, ob sie wahr oder falsch ist:
\par\smallskip\noindent \begingroup\renewcommand{\arraystretch}{1.5}\begin{tabular}{|p{8.0cm}|c|c|}\hline Aussage & wahr & falsch \\ \hline Der Punkt $(-1\,|\,7)$ liegt auf $p$. & $\square$ & $\square$ \\ \hline $p$ ist die Normalparabel, um 2 nach rechts und um 2 nach unten verschoben. & $\square$ & $\square$ \\ \hline $p$ schneidet die y-Achse im Punkt $(2\,|\,0)$ & $\square$ & $\square$ \\ \hline\end{tabular}\endgroup\par
\fusshilfe{\mbox{37:~wahr; falsch}}
\end{pfaufg}
\pfgruppe{}
\begin{pfaufg}{P10 ’14}{\llap{\textasteriskcentered\,}38.}{}
Im Koordinatensystem sind die Parabel $p$ mit $p(x) = -x^2$ und die Gerade $g$ mit $g(x) = 2x - 3$ gezeichnet. Untersuche mit einer Rechnung, ob $S(-3\,|\,-9)$ ein gemeinsamer Punkt von $p$ und $g$ ist.
\par\smallskip\noindent \begin{tikzpicture}\begin{axis}[xmin=-6,xmax=6,ymin=-6,ymax=6,axis lines=middle,tick label style={font=\scriptsize},clip=true,/pgf/number format/use comma,axis line style={-{Stealth[length=2mm]}},every axis plot/.append style={line width=0.9pt},x=0.55cm,y=0.55cm,xtick distance=1,ytick distance=1,enlargelimits=false,grid=both,grid style={mbgitter,line width=0.3pt},major grid style={mbgitter},xlabel={$x$},ylabel={$y$},x label style={anchor=west},y label style={anchor=south}]\addplot[black,domain=-6:6,samples=120] {-(x^2)};\node[font=\small,fill=white,inner sep=1pt,anchor=south west] at (axis cs:2.136,-4.5625) {$p$};\addplot[black,domain=-6:6,samples=120] {2*x-3};\node[font=\small,fill=white,inner sep=1pt,anchor=south west] at (axis cs:3.72,4.44) {$g$};\end{axis}\end{tikzpicture}\par
\rechenplatz{2}
\fusshilfe{\mbox{38:~ja, S($-$3 | $-$9) ist Schnittpunkt}}
\end{pfaufg}

\pfstufekopf[sparabelskizzieren]{Parabel skizzieren}{in 1 der letzten 5 Prüfungen}
\pfgruppe{}
\begin{pfaufg}{P10 ’24}{39.}{}
Die Gerade $f$ ist der Graph von $f(x) = 4x + 1$. Die Parabel $p$ ist der Graph von $p(x) = x^2 - 4$. Skizziere die Parabel $p$ in das Koordinatensystem. Gib ihren Scheitelpunkt an: $S(\ \ |\ \ )$.
\par\smallskip\noindent \begin{tikzpicture}\begin{axis}[xmin=-4,xmax=4,ymin=-5,ymax=8,axis lines=middle,tick label style={font=\scriptsize},clip=true,/pgf/number format/use comma,axis line style={-{Stealth[length=2mm]}},every axis plot/.append style={line width=0.9pt},x=0.42cm,y=0.42cm,xtick distance=1,ytick distance=1,enlargelimits=false,grid=both,grid style={mbgitter,line width=0.3pt},major grid style={mbgitter},xlabel={$x$},ylabel={$y$},x label style={anchor=west},y label style={anchor=south}]\end{axis}\end{tikzpicture}\par
\rechenplatz{1}
\fusshilfe{\mbox{39:~S(0|$-$4)}}
\end{pfaufg}

\pfstufekopf[slagezweierparabelnoh]{Lage zweier Parabeln ohne Rechnung begründen}{in 1 der letzten 5 Prüfungen}
\pfgruppe{}
\begin{pfaufg}{P10 ’14}{\llap{\textasteriskcentered\,}40.}{}
Im Koordinatensystem sind die Parabel $p$ mit $p(x) = -x^2$ und die Gerade $g$ mit $g(x) = 2x - 3$ gezeichnet. Eine weitere Gerade $f$ soll die Parabel $p$ nirgends treffen. Gib eine Gleichung für eine solche Gerade an.
\par\smallskip\noindent \begin{tikzpicture}\begin{axis}[xmin=-6,xmax=6,ymin=-6,ymax=6,axis lines=middle,tick label style={font=\scriptsize},clip=true,/pgf/number format/use comma,axis line style={-{Stealth[length=2mm]}},every axis plot/.append style={line width=0.9pt},x=0.55cm,y=0.55cm,xtick distance=1,ytick distance=1,enlargelimits=false,grid=both,grid style={mbgitter,line width=0.3pt},major grid style={mbgitter},xlabel={$x$},ylabel={$y$},x label style={anchor=west},y label style={anchor=south}]\addplot[black,domain=-6:6,samples=120] {-(x^2)};\node[font=\small,fill=white,inner sep=1pt,anchor=south west] at (axis cs:2.136,-4.5625) {$p$};\addplot[black,domain=-6:6,samples=120] {2*x-3};\node[font=\small,fill=white,inner sep=1pt,anchor=south west] at (axis cs:3.72,4.44) {$g$};\end{axis}\end{tikzpicture}\par
\rechenplatz{1}
\fusshilfe{\mbox{40:~z. B. f(x) = 1}}
\end{pfaufg}
\pfstufekopf[slsungprfen]{Lösung prüfen}{in 2 der letzten 5 Prüfungen}
\pfgruppe{}
\begin{pfaufg}{}{41.}{}
Kreuze an, wie viele Lösungen die Gleichung $(x - 3)^2 = 25$ hat.
\pfkreuzzeile{zwei}\pfkreuzzeile{eine}\pfkreuzzeile{keine}
\fusshilfe{\mbox{41:~rechts positiv}}
\end{pfaufg}
\begin{pfaufg}{}{42.}{}
Der Satz vom Nullprodukt sagt: Ein Produkt ist null, wenn ein Faktor null ist. Kreuze an, ob der Satz für die Gleichung $(x - 2) \cdot (x + 7) = 0$ gilt.
\pfkreuzzeile{gilt}\pfkreuzzeile{gilt nicht}
\fusshilfe{\mbox{42:~rechts steht null}}
\end{pfaufg}
\pfgruppe{}
\begin{pfaufg}{P10 ’25}{43.}{}
Genau einer der Werte löst die Gleichung $x(x+5) = -6$. Kreuze ihn an:
\pfkreuzzeile{$x = 2$}\pfkreuzzeile{$x = 3$}\pfkreuzzeile{$x = -2$}\pfkreuzzeile{$x = -6$}
\fusshilfe{\mbox{43:~x = $-$2}}
\end{pfaufg}
\pfstufekopf[sxzugegebenemy]{x zu gegebenem y}{in 1 der letzten 5 Prüfungen}
\pfgruppe{}
\begin{pfaufg}{}{44.}{}
Gegeben ist $f(x) = x^2 + x + 1$. Für welche $x$ ist $f(x) = 7$?
\par\noindent x\_1 = \leerfeld , x\_2 = \leerfeld \par
\fusshilfe{\mbox{44:~$x_1 = 2$, $x_2 = -3$}}
\end{pfaufg}
\pfgruppe{}
\begin{pfaufg}{P10 ’23}{\llap{\textasteriskcentered\,}45.}{}
Die Parabel $p$ hat die Gleichung $p(x) = -x^2 - 2x + 5$. Bestimme die x-Koordinaten aller Punkte von $p$, deren y-Koordinate $-10$ ist.
\rechenplatz{3}
\fusshilfe{\mbox{45:~x$_1$ = $-$5, x$_2$ = 3}}
\end{pfaufg}
\par\Needspace*{0.55\textheight}\pfstufekopf[serkennen1]{Was ist gesucht?}{}
\begin{pfaufg}{}{46.}{}
Was ist gesucht? Kreuze an. Rechne nicht.
\par\smallskip\noindent \begingroup\renewcommand{\arraystretch}{1.4}\begin{tabular}{|>{\raggedright\arraybackslash}p{8.3cm}|>{\centering\arraybackslash\hyphenpenalty=0}p{1.55cm}|>{\centering\arraybackslash\hyphenpenalty=0}p{1.55cm}|}\hline  & {\scriptsize\hspace{0pt}Nullstellen berechnen} & {\scriptsize\hspace{0pt}x zu gegebenem y} \\ \hline Wo hat y = (x + 3)² $-$ 2 seine Nullstellen? & $\square$ & $\square$ \\ \hline Welche Punkte der Parabel y = $-$x² $-$ 2x + 5 haben die y-Koordinate $-$10? & $\square$ & $\square$ \\ \hline Berechne die Nullstellen von y = 2x² + 8x + 6. & $\square$ & $\square$ \\ \hline Die Punkte P$_1$(x$_1$ | 23) und P$_2$(x$_2$ | 23) liegen auf y = x² $-$ 2. Bestimme x$_1$ und x$_2$. & $\square$ & $\square$ \\ \hline In welchen Punkten schneidet die Parabel y = x² $-$ 6x + 7 die x-Achse? & $\square$ & $\square$ \\ \hline\end{tabular}\endgroup\par
\end{pfaufg}
\pfpruefstein{P10 ’26}
\begin{pfaufg}{}{}{}
Die verschobene Normalparabel $p$ ist der Graph von $p(x) = (x-2)^2 - 2$ und ist im Koordinatensystem gezeichnet.
\par\smallskip\noindent \begin{tikzpicture}\begin{axis}[xmin=-2,xmax=5,ymin=-2,ymax=5,axis lines=middle,tick label style={font=\scriptsize},clip=true,/pgf/number format/use comma,axis line style={-{Stealth[length=2mm]}},every axis plot/.append style={line width=0.9pt},x=0.55cm,y=0.55cm,xtick distance=1,ytick distance=1,enlargelimits=false,grid=both,grid style={mbgitter,line width=0.3pt},major grid style={mbgitter},xlabel={$x$},ylabel={$y$},x label style={anchor=west},y label style={anchor=south}]\addplot[black,domain=-2:5,samples=120] {((x-2)^2)-2};\node[font=\small,fill=white,inner sep=1pt,anchor=south west] at (axis cs:4.44,3.9536) {$p$};\end{axis}\end{tikzpicture}\par
\end{pfaufg}
\begin{pfaufg}{}{*a)}{2 BE}
Die lineare Funktion $f$ hat die Gleichung $f(x) = -2x + 2$. Zeichne ihren Graphen in das Koordinatensystem.
\rechenplatz{3}
\fusshilfe{\mbox{Prüfstein a):~Gerade durch (0|2) und (1|0), fallend}}
\end{pfaufg}
\begin{pfaufg}{}{b)}{2 BE}
Prüfe mit einer Rechnung, ob der Punkt $P(-4\,|\,10)$ auf dem Graphen von $f(x) = -2x + 2$ liegt.
\rechenplatz{3}
\fusshilfe{\mbox{Prüfstein b):~Ja, f($-$4) = $-$2 · ($-$4) + 2 = 10}}
\end{pfaufg}
\begin{pfaufg}{}{c)}{1 BE}
Gib den Scheitelpunkt von $p$ an: $S(\ \ |\ \ )$.
\rechenplatz{3}
\fusshilfe{\mbox{Prüfstein c):~S(2|$-$2)}}
\end{pfaufg}
\begin{pfaufg}{}{*d)}{3 BE}
Eine zweite Parabel $q$ hat die Gleichung $q(x) = -(x-2)^2 - 2$. Gib an, wie viele Punkte $p$ und $q$ gemeinsam haben. Begründe ohne zu rechnen.
\rechenplatz{3}
\fusshilfe{\mbox{Prüfstein d):~genau ein gemeinsamer Punkt S(2|$-$2)}}
\end{pfaufg}
\end{document}